\section{Methodology}
\label{sec:methodology}

\subsection{Notation}
\label{sec:notation}

Let $G = (V, E, w)$ be the census-tract adjacency graph for a state, where
$V$ is the set of census tracts, $E$ is the set of shared-boundary edges, and
$w: E \to \mathbb{R}_{>0}$ assigns each edge the TIGER boundary length between
adjacent tracts. Let $\text{pop}: V \to \mathbb{Z}_{>0}$ be the 2020 census
population of each tract. Let $n$ be the state's apportioned congressional seat
count.

\subsection{The Split Prescription}
\label{sec:factorization}

At each level of the PFR tree, a region with $k$ target districts is split
according to the following rule with threshold $\kappa = 3$:

\begin{definition}[Split prescription]
  \label{def:split}
  For $k \geq 1$, $\text{split}(k)$ is:
  \begin{enumerate}
    \item \textbf{Direct} ($k \leq \kappa$): split into $k$ equal parts
          (sub-regions each with target~1).
    \item \textbf{Uniform-composite} ($k > \kappa$, $k$ composite): let $p$ be
          the \emph{largest} prime factor of $k$. Split into $p$ equal parts,
          each with target $k/p$.
    \item \textbf{Binary fallback} ($k > \kappa$, $k$ prime): split into 2 parts
          with population targets $\lfloor k/2 \rfloor$ and $\lceil k/2 \rceil$.
  \end{enumerate}
\end{definition}

Examples: $\text{split}(8) = (2, \text{sub}=4)$; $\text{split}(9) = (3, \text{sub}=3)$;
$\text{split}(17) = (2, \text{sub}=(8,9))$; $\text{split}(34) = (17, \text{sub}=2)$;
$\text{split}(51) = (17, \text{sub}=3)$.

\paragraph{Why largest prime first?} The reuse property requires that two seat
counts $n$ and $n'$ sharing a common factor $p$ produce the same top-level
$p$-way partition. With largest-prime-first, $n=34=2\times 17$ and
$n'=51=3\times 17$ both prescribe a 17-way primary split, sharing the same
17-district top-level partition of any given state. Smallest-prime-first would
instead bisect the state first (different halves for $n=34$ vs.\ thirds for
$n'=51$), destroying reuse.

\paragraph{Why binary fallback for large primes?} A direct $k$-way METIS
cut for large prime $k$ (e.g.\ $k=17$) is computationally feasible but
produces no reuse with neighbouring seat counts and is harder to balance than
iterated bisections. The binary fallback $17 \to (8, 9)$ naturally decomposes
into the highly-reused sub-trees $k=8=2^3$ (three binary levels) and $k=9=3^2$
(two ternary levels).

\subsection{The PFR Algorithm}
\label{sec:pfr}

\begin{definition}[Prime-Factored Redistricting]
  Given graph $G$ and seat count $n$, the PFR map $M_n(G)$ is constructed by
  the recursive procedure $\text{PFR}(G, n)$:
  \begin{enumerate}
    \item If $n = 1$: return $\{G\}$ (one district = whole region).
    \item Compute $\text{split}(n) = (p, k_1, \ldots, k_p)$ where $p$ is the
          number of parts and $k_i$ is the target district count for part $i$
          (Definition~\ref{def:split}).
    \item Partition $G$ into $p$ sub-regions $R_1, \ldots, R_p$ with target
          population fractions $k_i / n$ each, minimising edge-cut weight.
    \item For each $i$: recursively compute $\text{PFR}(R_i, k_i)$.
    \item Return the union of all recursive results.
  \end{enumerate}
\end{definition}

The $p$-way partition in step~3 is computed by METIS using population vertex
weights and TIGER boundary-length edge weights. For the uniform case ($k_1 =
\cdots = k_p = n/p$) METIS uses equal target fractions; for the binary fallback
case ($k_1 = \lfloor n/2 \rfloor$, $k_2 = \lceil n/2 \rceil$) it uses
proportional target fractions.

\subsection{The Reuse Property}
\label{sec:reuse}

The key property of PFR is that splits are canonical:

\begin{definition}[Canonical $p$-split]
  For a state with graph $G$ and a prime $p$, the \emph{canonical $p$-split}
  $S_p(G)$ is the minimum $p$-way partition of $G$. For a sub-region
  $R \subseteq G$, $S_p(R)$ is the minimum $p$-way partition of the induced
  subgraph $G[R]$.
\end{definition}

\begin{theorem}[Reuse]
  \label{thm:reuse}
  Let $n$ and $n'$ be composite integers with the same largest prime factor $p$,
  i.e.\ $n = p \cdot m$ and $n' = p \cdot m'$ with $p > \max(\text{prime
  factors of } m)$ and $p > \max(\text{prime factors of } m')$. Then
  $M_n(G)$ and $M_{n'}(G)$ share the same top-level $p$-way partition
  $S_p(G)$.
\end{theorem}

\begin{proof}
  Since $p$ is the largest prime factor of both $n$ and $n'$, the split
  prescription (Definition~\ref{def:split}, case 2) applies $p$ at level 1
  for both. The minimum $p$-way partition of $G$ is uniquely determined
  (assuming no ties) by $G$, $p$, and the METIS seed. Hence both maps share
  $S_p(G)$ at level 1. \qed
\end{proof}

\begin{corollary}
  $M_{34}(G)$ and $M_{51}(G)$ share the same 17-way top-level partition of $G$,
  since $34 = 2 \times 17$ and $51 = 3 \times 17$ both have largest prime 17.
  After the shared first level, $M_{34}$ bisects each of the 17 sub-regions
  while $M_{51}$ trisects them.
\end{corollary}

More generally, any two seat counts $n, n'$ that agree on their first $k$ terms
of $F(n)$ share the first $k$ levels of their factorization tree.

\subsection{Factorization Tree Structure}
\label{sec:tree}

The factorization tree $T_n$ for a state has:
\begin{itemize}
  \item \textbf{Root}: the whole state.
  \item \textbf{Level $\ell$ nodes}: regions produced by applying $q_1, \ldots,
        q_\ell$ to the root. Each level-$\ell$ node has $q_\ell$ children.
  \item \textbf{Leaves}: the $n$ districts.
  \item \textbf{Depth}: $\Omega(n)$ (number of prime factors with multiplicity).
\end{itemize}

For $n = 2^k$ (e.g., Wisconsin: $n = 8 = 2^3$), the tree is a complete binary
tree of depth~$k$ --- exactly the recursive bisection algorithm of
\citet{dellaLibera2026mec}.

For $n$ prime (e.g., Pennsylvania: $n = 17$), the tree has depth~1: a single
$n$-way cut of the whole state. No reuse is possible across different prime seat
counts.

\subsection{Handling Reapportionment}
\label{sec:reapportionment}

When $n$ changes to $n' = n \pm 1$ after a decennial census, the required
recomputation depends on the first position where $F(n)$ and $F(n')$ differ.

\begin{proposition}[Recomputation depth]
  Let $F(n) = [q_1, \ldots, q_d]$ and $F(n') = [q_1', \ldots, q_{d'}']$.
  Let $k^* = \min\{k : q_k \neq q_k'\}$ (or $k^* = 1$ if $q_1 \neq q_1'$).
  Then PFR reuses levels $1, \ldots, k^*-1$ and recomputes levels
  $k^*, \ldots, \max(d, d')$.
\end{proposition}

The worst case is $k^* = 1$: the largest prime factor changes, requiring a
full top-level recomputation. This occurs for the $51 \to 52$ transition:
$51 = 3 \times 17$ (largest prime 17) vs.\ $52 = 4 \times 13$ (largest prime 13).
No geographic sub-region is shared between the two maps.

\subsection{PFR-Smooth: Minimising Reapportionment Disruption}
\label{sec:smooth}

When the prime factorization changes structurally, PFR-smooth finds the
largest common prefix of $F(n)$ and $F(n')$ and reuses it, completing the
remaining levels from the best available cached splits. Specifically, for
$n \to n+1$:
\begin{enumerate}
  \item Find the longest common prefix of $F(n)$ and $F(n+1)$.
  \item Reuse those levels' cached splits.
  \item Re-solve only the subtrees that changed.
\end{enumerate}
This minimises the number of census tracts that change district assignment
across reapportionments.
